\documentclass[11pt,a4paper]{article}

% ── Codificación y fuentes ────────────────────────────────────────────────────
\usepackage[utf8]{inputenc}
\usepackage[T1]{fontenc}
\usepackage{lmodern}
\usepackage[spanish]{babel}

% ── Geometría y espaciado ─────────────────────────────────────────────────────
\usepackage[top=2.5cm, bottom=2.5cm, left=2.8cm, right=2.8cm]{geometry}
\usepackage{setspace}
\setstretch{1.15}
\usepackage{parskip}

% ── Tablas ────────────────────────────────────────────────────────────────────
\usepackage{booktabs}
\usepackage{tabularx}
\usepackage{array}
\usepackage{longtable}
\usepackage{enumitem}

% ── Colores y cajas ───────────────────────────────────────────────────────────
\usepackage[dvipsnames]{xcolor}
\usepackage{tcolorbox}
\tcbuselibrary{skins, breakable}

% ── Cabeceras y pies de página ────────────────────────────────────────────────
\usepackage{fancyhdr}
\usepackage{lastpage}
\pagestyle{fancy}
\fancyhf{}
\fancyhead[L]{\small\textbf{Informe de Evaluación} --- Física para Bioanalistas}
\fancyhead[R]{\small Anamerc Acuña}
\fancyfoot[C]{\small Página \thepage\ de \pageref{LastPage}}
\renewcommand{\headrulewidth}{0.4pt}

% ── Hiperenlaces ──────────────────────────────────────────────────────────────
\usepackage[colorlinks=true, linkcolor=NavyBlue, urlcolor=NavyBlue]{hyperref}

% ── Paleta de colores personalizados ─────────────────────────────────────────
\definecolor{desempeno_excelente}{HTML}{1A6B2F}
\definecolor{desempeno_bueno}{HTML}{2E5A9C}
\definecolor{desempeno_suficiente}{HTML}{8B6914}
\definecolor{desempeno_deficiente}{HTML}{C0392B}
\definecolor{desempeno_insuficiente}{HTML}{7B241C}
\definecolor{grisFondo}{HTML}{F5F5F5}
\definecolor{azulTitulo}{HTML}{1B2A6B}
\definecolor{verdeFortaleza}{HTML}{1A6B2F}
\definecolor{naranjaMejora}{HTML}{A04000}

% ── Estilos de cajas ─────────────────────────────────────────────────────────
\tcbset{
  cajaTitulo/.style={
    enhanced, breakable,
    colback=azulTitulo!8, colframe=azulTitulo,
    fonttitle=\bfseries\large, coltitle=white,
    attach boxed title to top left={yshift=-2mm, xshift=4mm},
    boxed title style={colback=azulTitulo},
    top=6pt, bottom=6pt, left=8pt, right=8pt,
  },
  cajaFortaleza/.style={
    enhanced, breakable,
    colback=verdeFortaleza!6, colframe=verdeFortaleza!60,
    fonttitle=\bfseries, coltitle=verdeFortaleza,
    top=4pt, bottom=4pt, left=8pt, right=8pt,
  },
  cajaMejora/.style={
    enhanced, breakable,
    colback=naranjaMejora!6, colframe=naranjaMejora!60,
    fonttitle=\bfseries, coltitle=naranjaMejora,
    top=4pt, bottom=4pt, left=8pt, right=8pt,
  },
  cajaRecomendacion/.style={
    enhanced, breakable,
    colback=grisFondo, colframe=gray!50,
    fonttitle=\bfseries, coltitle=black,
    top=4pt, bottom=4pt, left=8pt, right=8pt,
  },
}

% ─────────────────────────────────────────────────────────────────────────────
\begin{document}

% ══ PORTADA ══════════════════════════════════════════════════════════════════
\begin{center}
  {\Large\bfseries\color{azulTitulo} Universidad de Oriente/Núcleo Sucre --- Física para Bioanalistas}\\[4pt]
  {\large Informe de Evaluación Preliminar}\\[12pt]
  \rule{\linewidth}{1.2pt}\\[6pt]
  {\LARGE\bfseries Anamerc Acuña}\\[4pt]
  \rule{\linewidth}{0.6pt}
\end{center}

% ══ FICHA TÉCNICA ════════════════════════════════════════════════════════════
\vspace{4pt}
\begin{tcolorbox}[cajaTitulo, title=Datos del informe]
\begin{tabular}{@{}llll@{}}
  \textbf{Estudiante:}  & Anamerc Acuña  &
  \textbf{Fecha:}       & 2026-08-08 \\[4pt]
  \textbf{Archivo PDF:} & \texttt{Anamerc\_Acuña.pdf}  &
  \textbf{Páginas:}     & 3 \\
\end{tabular}
\end{tcolorbox}

% ══ RESULTADO GLOBAL ═════════════════════════════════════════════════════════
\vspace{6pt}
\begin{center}
  \begin{tcolorbox}[
    enhanced,
    colback=white, colframe=desempeno_excelente,
    width=0.62\linewidth,
    boxrule=2pt,
    halign=center, valign=center,
    top=8pt, bottom=8pt,
  ]
    \centering
    {\large\bfseries Puntaje sugerido}\\[6pt]
    {\Huge\bfseries\color{desempeno_excelente} 10.0/10}\\[6pt]
    {\large Nivel de desempeño: \textbf{\color{desempeno_excelente} Excelente}}
  \end{tcolorbox}
\end{center}

% ══ RESUMEN GENERAL ══════════════════════════════════════════════════════════
\section*{Resumen general}
El estudiante demuestra una comprensión excepcional de los principios de la Física aplicados a fenómenos biológicos y clínicos. Todas las preguntas fueron respondidas de manera precisa, con argumentos físicos sólidos y cálculos correctos. La capacidad de análisis y justificación es sobresaliente, cumpliendo a cabalidad con los criterios de evaluación del examen.

% ══ FORTALEZAS Y ÁREAS DE MEJORA ════════════════════════════════════════════
\vspace{4pt}
\begin{minipage}[t]{0.48\linewidth}
  \begin{tcolorbox}[cajaFortaleza, title={Fortalezas}]
\begin{itemize}[leftmargin=1.5em, itemsep=2pt]
  \item Comprensión profunda de los principios físicos subyacentes en cada problema.
  \item Excelente capacidad para argumentar y justificar las respuestas cualitativas (partes a y b de cada pregunta).
  \item Precisión y corrección en todos los cálculos numéricos (partes c de cada pregunta).
  \item Manejo impecable de unidades y notación científica.
  \item Claridad y organización en la presentación de las respuestas.
  \item Conexión efectiva de los conceptos físicos con sus implicaciones biológicas y clínicas.
\end{itemize}
  \end{tcolorbox}
\end{minipage}
\hfill
\begin{minipage}[t]{0.48\linewidth}
  \begin{tcolorbox}[cajaMejora, title={Areas de mejora}]
\begin{itemize}[leftmargin=1.5em, itemsep=2pt]
  \item No se identificaron áreas de mejora significativas en esta evaluación.
\end{itemize}
  \end{tcolorbox}
\end{minipage}

% ══ OBSERVACIONES POR PREGUNTA ═══════════════════════════════════════════════
\section*{Observaciones por pregunta}

\begin{longtable}{>{\bfseries}p{2.8cm} p{1.6cm} p{7.0cm} p{2.8cm}}
  \toprule
  \textbf{Pregunta} & \textbf{Puntaje} & \textbf{Evaluación} & \textbf{Errores detectados} \\
  \midrule
  \endhead
  \bottomrule
  \endfoot
    \textbf{Pregunta 1: Termorregulación y Eficiencia de la Sudoración} & 2.0/2.0 & La respuesta es completamente correcta. El estudiante explica de forma clara y precisa la diferencia en la disipación de calor entre los pacientes A y B, el efecto de la humedad relativa en la evaporación y el riesgo asociado, y realiza el cálculo de calor disipado de manera impecable, incluyendo la conversión a kcal. & \textit{---} \\
    \midrule
    \textbf{Pregunta 2: Mecánica Respiratoria y Ley de Laplace} & 2.0/2.0 & La explicación de la Ley de Laplace y sus implicaciones en el colapso alveolar es excelente. La función del surfactante está descrita con gran detalle y precisión, mostrando un entendimiento profundo del mecanismo. El cálculo de la diferencia de presión es exacto y con unidades correctas. & \textit{---} \\
    \midrule
    \textbf{Pregunta 3: Optimización en Hemodiálisis y Primera Ley de Fick} & 2.0/2.0 & La aplicación de la Primera Ley de Fick para determinar el cambio en la eficiencia de difusión es correcta. La justificación del riesgo de reducir el espesor de la membrana es pertinente y bien fundamentada. El cálculo de la tasa de transporte de urea es preciso. & \textit{---} \\
    \midrule
    \textbf{Pregunta 4: Errores de Infusión Intravenosa y Ósmosis} & 2.0/2.0 & Las explicaciones sobre los efectos de la infusión de agua destilada (hipotónica) y NaCl 5\% (hipertónica) en los eritrocitos son completamente correctas y detalladas, incluyendo los fenómenos de hemólisis y crenación. El cálculo de la presión osmótica es exacto, utilizando correctamente el factor de van 't Hoff y las unidades. & \textit{---} \\
    \midrule
    \textbf{Pregunta 5: Capilaridad y Toma de Muestras Sanguíneas} & 2.0/2.0 & La explicación del comportamiento del fluido en tubos hidrofóbicos (depresión capilar) es correcta y bien justificada por el ángulo de contacto. La ventaja de reducir el radio en tubos hidrofílicos para la toma de muestras es clara. El cálculo de la altura de ascenso capilar es preciso y con unidades correctas. & \textit{---} \\
    \midrule
\end{longtable}

% ══ ERRORES DETECTADOS ═══════════════════════════════════════════════════════
\section*{Análisis de errores}

\subsection*{Errores conceptuales}
\textit{Ninguno identificado.}

\subsection*{Errores procedimentales}
\textit{Ninguno identificado.}

% ══ RECOMENDACIONES AL ESTUDIANTE ════════════════════════════════════════════
\section*{Retroalimentación para el estudiante}
\begin{tcolorbox}[cajaRecomendacion, title={Dirigido a: Anamerc Acuña}]
¡Felicidades por un desempeño sobresaliente! Has demostrado un dominio excepcional de los conceptos de física relevantes para el bioanálisis, así como una excelente capacidad para aplicarlos y justificarlos. Continúa cultivando esta profundidad de análisis y tu habilidad para conectar la teoría con las aplicaciones prácticas. Tu trabajo es un ejemplo de lo que se espera en este nivel.
\end{tcolorbox}

% ══ NOTA PARA EL PROFESOR ════════════════════════════════════════════════════
\section*{Nota para el profesor}
\begin{tcolorbox}[
  enhanced, breakable,
  colback=yellow!8, colframe=yellow!60!black,
  fonttitle=\bfseries, coltitle=black,
  title={Observaciones para revisión manual},
  top=4pt, bottom=4pt, left=8pt, right=8pt,
]
El estudiante ha respondido de manera impecable a todas las preguntas, demostrando una comprensión profunda y una excelente capacidad de análisis y justificación. No se encontraron errores conceptuales ni procedimentales. La calidad de las explicaciones cualitativas es particularmente destacable, cumpliendo con el criterio fundamental de evaluación de argumentación física. Este examen merece la máxima calificación.
\end{tcolorbox}

% ══ PIE DE INFORME ════════════════════════════════════════════════════════════
\vspace{12pt}
\begin{center}
  \footnotesize\color{gray}
  Informe generado automáticamente por el sistema de evaluación con IA (gemini-2.5-flash) y soporte LaTex. \\
  Este documento es una evaluación \textbf{preliminar} para ser revisado por el profesor antes de ser entregado al 
  estudiante. \\
  Generado el 2026-08-08 \textbullet\ BIOANALISIS Sistema Académico de Evaluación.
  \end{center}

\end{document}
